\section{Valid Parenthesis String}
\label{sec:valid-parenthesis-string}

\problemstatement{We are given a string $s$ consisting only of the
characters \texttt{'('}, \texttt{')'}, and \texttt{'*'}. The character
\texttt{'*'} is a wildcard: each occurrence of it may be treated, independently,
as either \texttt{'('}, or \texttt{')'}, or the empty string. We must decide
whether there is \emph{some} way to resolve every \texttt{'*'} so that the
resulting string is a valid parenthesis sequence (every prefix has at least
as many \texttt{'('} as \texttt{')'}, and the total counts of \texttt{'('}
and \texttt{')'} are equal).}

\begin{patternbox}
The phrase \emph{``each wildcard may independently be one of several
things; does \textbf{some} assignment make the whole thing valid''} looks,
at first glance, like it demands trying all $3^k$ assignments of the $k$
wildcards -- an exponential brute force. But the underlying validity
condition for parentheses is itself expressible as a simple running balance
(``open count never goes negative, and ends at zero''), and whenever a
validity condition can be tracked by a \emph{running numeric range} rather
than an exact value, that is a strong signal that we can collapse all
exponentially many choices into tracking just the minimum and maximum
possible value of that running quantity at each step -- a classic
``track an interval of reachable states'' greedy, no recursion or
backtracking required.
\end{patternbox}

\subsection*{Understanding the problem carefully}

For a string with only \texttt{'('} and \texttt{')'}, validity is checked by
maintaining a single counter, the \emph{open balance}: add $1$ for
\texttt{'('}, subtract $1$ for \texttt{')'}, and the string is valid exactly
when the balance never goes negative and ends at exactly $0$.

With \texttt{'*'} in the mix, the open balance after processing a prefix is
no longer a single number -- it is a \emph{range} of possible numbers,
depending on how each \texttt{'*'} so far was resolved. Let
$\textit{lo}$ be the smallest possible balance consistent with some
resolution of the wildcards seen so far, and $\textit{hi}$ be the largest
possible balance. We update this range character by character:

\begin{itemize}
  \item On \texttt{'('}: both the smallest and largest possible balance
        increase by $1$, so $\textit{lo} \mathrel{+}= 1$,
        $\textit{hi} \mathrel{+}= 1$.
  \item On \texttt{')'}: both decrease by $1$, so
        $\textit{lo} \mathrel{-}= 1$, $\textit{hi} \mathrel{-}= 1$.
  \item On \texttt{'*'}: in the most pessimistic resolution (treat it as
        \texttt{')'}), the balance decreases by $1$, so
        $\textit{lo} \mathrel{-}= 1$. In the most optimistic resolution
        (treat it as \texttt{'('}), the balance increases by $1$, so
        $\textit{hi} \mathrel{+}= 1$. (Treating it as empty would leave the
        balance unchanged, which always lies between these two extremes, so
        it never needs separate tracking.)
\end{itemize}

\begin{ideabox}[Greedy Choice]
At every position, maintain only the \emph{interval} $[\textit{lo},
\textit{hi}]$ of balances that are achievable by \emph{some} resolution of
the wildcards seen so far -- not every individual achievable value, just the
two extremes. Whenever $\textit{hi}$ drops below $0$, no resolution can
recover (the balance has gone irrecoverably negative even in the best case),
so we can answer \texttt{false} immediately. Whenever $\textit{lo}$ drops
below $0$, clamp it to $0$, because a balance can never validly be negative,
so any resolution that would have produced a negative running balance is
simply discarded from consideration -- it is never the right choice for an
eventual valid string, so we do not need to track it.
\end{ideabox}

\subsection*{Why tracking only the interval is enough}

It might seem like we are throwing away information by tracking only
$[\textit{lo}, \textit{hi}]$ rather than the full set of achievable balances.
But notice that the set of achievable balances after processing a prefix is
always a contiguous range of integers with the \emph{same parity step of
$1$} apart (because every character changes the balance by exactly $-1$,
$0$, or $+1$, and a \texttt{'*'} contributes all three options, which fills
in every integer between the all-\texttt{')'} extreme and the all-\texttt{'('}
extreme). So the reachable set truly is an interval, and tracking its
endpoints loses nothing. At the very end, the string can be made valid if
and only if $0$ lies in the final interval -- but since we always clamp
$\textit{lo}$ to be at least $0$ throughout, checking that the final
$\textit{lo} = 0$ is equivalent to checking that $0$ is reachable, because if
$0$ were reachable, clamping would have kept $\textit{lo}$ pinned at (or
below, then clamped to) $0$ rather than drifting strictly positive.

\subsection*{Algorithm}

\begin{enumerate}
  \item Initialize $\textit{lo} \gets 0$, $\textit{hi} \gets 0$.
  \item For each character $c$ in $s$:
  \begin{enumerate}
    \item If $c = $ \texttt{'('}: $\textit{lo} \gets \textit{lo}+1$,
          $\textit{hi} \gets \textit{hi}+1$.
    \item If $c = $ \texttt{')'}: $\textit{lo} \gets \textit{lo}-1$,
          $\textit{hi} \gets \textit{hi}-1$.
    \item If $c = $ \texttt{'*'}: $\textit{lo} \gets \textit{lo}-1$,
          $\textit{hi} \gets \textit{hi}+1$.
    \item If $\textit{hi} < 0$: return \texttt{false} immediately.
    \item If $\textit{lo} < 0$: clamp $\textit{lo} \gets 0$.
  \end{enumerate}
  \item After processing the whole string, return
        $(\textit{lo} = 0)$.
\end{enumerate}

\begin{algorithm}[H]
\caption{Valid Parenthesis String}
\begin{algorithmic}[1]
\Procedure{CheckValidString}{$s$}
  \State $\textit{lo} \gets 0$, $\textit{hi} \gets 0$
  \For{each character $c$ in $s$}
    \If{$c = $ '('}
      \State $\textit{lo} \gets \textit{lo} + 1$, $\textit{hi} \gets \textit{hi} + 1$
    \ElsIf{$c = $ ')'}
      \State $\textit{lo} \gets \textit{lo} - 1$, $\textit{hi} \gets \textit{hi} - 1$
    \Else \Comment{$c = $ '*'}
      \State $\textit{lo} \gets \textit{lo} - 1$, $\textit{hi} \gets \textit{hi} + 1$
    \EndIf
    \If{$\textit{hi} < 0$}
      \State \Return \textbf{false}
    \EndIf
    \If{$\textit{lo} < 0$}
      \State $\textit{lo} \gets 0$
    \EndIf
  \EndFor
  \State \Return $(\textit{lo} = 0)$
\EndProcedure
\end{algorithmic}
\end{algorithm}

\subsection*{Dry run}

\dryrun{Take $s = \texttt{"(*))"}$.}

Start with $\textit{lo}=0$, $\textit{hi}=0$.

\begin{itemize}
  \item $c = $ \texttt{'('}: $\textit{lo}=1$, $\textit{hi}=1$.
  \item $c = $ \texttt{'*'}: $\textit{lo}=0$, $\textit{hi}=2$.
  \item $c = $ \texttt{')'}: $\textit{lo}=-1 \to$ clamp to $0$,
        $\textit{hi}=1$.
  \item $c = $ \texttt{')'}: $\textit{lo}=-1 \to$ clamp to $0$,
        $\textit{hi}=0$.
\end{itemize}

At no point did $\textit{hi}$ go negative, and the final $\textit{lo} = 0$,
so the answer is \texttt{true}. We can confirm this directly: resolving the
\texttt{'*'} as \texttt{'('} turns $s$ into \texttt{"(())"}, whose balances
are $1,2,1,0$ -- never negative, ending at $0$ -- so this is indeed a valid
resolution, and the algorithm's \texttt{true} verdict is justified.

Now take $s = \texttt{"(*))("}$, which appends one extra unmatched
\texttt{'('} to the previous string.

\begin{itemize}
  \item As above, after \texttt{"(*))"}: $\textit{lo}=0$, $\textit{hi}=0$.
  \item $c = $ \texttt{'('}: $\textit{lo}=1$, $\textit{hi}=1$.
\end{itemize}

Final $\textit{lo} = 1 \neq 0$, so the answer is \texttt{false} -- correctly,
since the trailing unmatched \texttt{'('} can never be closed by anything
after it.

\subsection*{C++ implementation}

\begin{lstlisting}
#include <bits/stdc++.h>
using namespace std;

bool checkValidString(string s) {
    int lo = 0, hi = 0;

    for (char c : s) {
        if (c == '(') {
            lo++;
            hi++;
        } else if (c == ')') {
            lo--;
            hi--;
        } else { // c == '*'
            lo--;
            hi++;
        }

        if (hi < 0) return false;
        if (lo < 0) lo = 0;
    }
    return lo == 0;
}

int main() {
    string s = "(*))";
    cout << (checkValidString(s) ? "true" : "false") << endl;
    return 0;
}
\end{lstlisting}

\begin{complexitybox}
\textbf{Time Complexity.} We scan the string once, doing constant work per
character, giving
\[
  O(n).
\]
\textbf{Space Complexity.} We maintain only the two integers \textit{lo} and
\textit{hi}, so the space complexity is $O(1)$.
\end{complexitybox}

\begin{takeawaybox}
Whenever a validity check can be expressed as a single running numeric
quantity that must stay within bounds, and the input introduces wildcards
with several possible numeric effects, try tracking the \emph{interval} of
reachable values for that quantity instead of branching over every
combination of wildcard resolutions. This turns an exponential
``try every assignment'' search into a single linear greedy sweep.
\end{takeawaybox}
